\section{Minecraft as a Construction Substrate}
\label{sec:minecraft-mechanics}

A central claim of this work is that building a quantum circuit inside Minecraft constitutes
a meaningfully different experience from assembling one in a schematic-based composer.
Substantiating that claim requires describing precisely what ``building'' means within
Minecraft itself, and identifying which of the game's own mechanics the mod reuses, extends,
or deliberately avoids.

\subsection{The voxel world and the timeline metaphor}

Minecraft's world is a three-dimensional grid of one-metre cubes, or \emph{blocks}; nearly
every action available to a player reduces to placing and removing blocks within this grid.
QuantumCraft maps a standard quantum-circuit diagram directly onto this grid: one qubit
occupies one row of qubit-line blocks running east, and a gate or marker block placed some
distance east of a qubit line's anchor represents an operation applied at that row's
corresponding time step, exactly as a gate drawn some distance along a horizontal wire in a
textbook circuit diagram represents an operation at that point in time. A two-qubit gate
such as CNOT is built from a pair of control and target marker blocks placed at the
\emph{same} east-offset on two different qubit rows, reproducing the vertical connecting line
a circuit diagram draws between a control dot and a target symbol, except that no physical
wire block is required between them: Section~\ref{sec:architecture} describes how the mod's
scanner recognizes this correlation purely from the two markers' shared time column, without
needing an explicit connector the way CircuitCraft's electrical topology
requires~\cite{picos2026circuitcraft}. This is a deliberate departure from CircuitCraft's own
graph/union--find topology model (in which any two mutually conductive faces are
electrically joined): a quantum circuit is conventionally drawn, taught, and reasoned about as
a fixed number of horizontal time-ordered lines, not as an arbitrary wiring graph, so the
in-world representation was chosen to match the diagram a student already encounters in a
textbook or lecture slide, rather than reusing CircuitCraft's own metaphor merely for
consistency between the two mods.

\subsection{Redstone: why it does not extend to qubits}

The base game already provides an entire circuit-like subsystem, redstone, built around
\emph{redstone dust} (a wire-like block carrying a signal), \emph{redstone torches} and other
sources, and a large family of blocks that read an incoming signal to trigger an effect. A
redstone signal is an integer strength from 0 to 15, propagated through a discrete
neighbor-notification event fired once per game tick; despite having no continuous quantity
(CircuitCraft's own paper discusses this same mechanic at greater length in the
classical-circuits context~\cite{picos2026circuitcraft}), redstone is widely used, both
informally and as a subject of academic interest~\cite{nebel2016mining}, as an entry point
into digital logic:
players construct logic gates, adders, and working memory cells entirely from redstone
components, and commonly describe and reason about redstone circuits in explicitly
electrical language~\cite{dezuanni2015redstone}. That established role, however, does not
extend to representing a qubit, and for a reason more fundamental than redstone's
already-discrete signal strength: a redstone signal, whatever its strength, is always a
single, locally determined, publicly readable integer, whereas a qubit's defining property is
that it exists in a superposition of two amplitudes that cannot both be read out at once, and
two qubits can become entangled such that neither one's state is even well defined
independently of the other. No quantization of a single scalar signal, however fine, can
represent a joint property that literally does not decompose into a value held separately by
each of two locations, which is precisely why QuantumCraft introduces an entirely new
representation (Section~\ref{sec:architecture}) rather than attempting to encode qubits into
redstone's existing signal-strength mechanic the way some ad hoc community techniques
approximate a continuous analog value with it.

\subsection{Block entities and on-demand execution}

Minecraft's modding interfaces expose the \emph{block entity} construct for a block that must
retain arbitrary mutable data, and, if the block type requests it, execute code once per game
tick. Every gate, qubit line, and control/target marker in this mod is implemented as a block
entity for the same reason CircuitCraft's components are~\cite{picos2026circuitcraft}: it is
the platform's existing mechanism for state tied to a specific position. Unlike
CircuitCraft, however, none of QuantumCraft's block entities register a per-tick callback at
all: a quantum circuit has no continuous transient to track between gate applications the way
a charging capacitor does, so there is nothing for a circuit to be doing between one probe use
and the next. A circuit is instead scanned and run once, on demand, exactly when a player
right-clicks the probe against a qubit line (Section~\ref{sec:architecture}), which is both
simpler to implement and a closer match to how a real quantum computer is actually used: a
circuit is prepared once and then executed as a single shot, not continuously re-evaluated in
the background.

\subsection{What the mod adds}

Given the preceding discussion, the mod's own contribution is a new, purpose-built spatial
notation for a quantum circuit -- qubit rows as parallel timelines, single-qubit gates as
blocks placed directly on those rows, and multi-qubit correlation expressed through shared
time-columns across rows rather than through any wired connection -- together with a
genuine statevector simulator to execute whatever circuit that notation describes
(Section~\ref{sec:quantum-simulator}). Redstone itself plays no role in this mod at all,
unlike in CircuitCraft, where it is reused deliberately for the one concept it is genuinely
well suited to represent (a boolean on/off switch); no comparably suitable role for redstone's
existing mechanics was identified for a quantum circuit, since even the mod's classical
special case (Section~\ref{sec:experiments}'s Toffoli-gate experiment) is more directly
expressed as boolean amplitudes at basis states $|0\rangle$ and $|1\rangle$ than through a
signal-strength integer.
